% BEGIN CR28_RADIAL_FULL
% Sección incorporable; requiere amsmath y amssymb.
% Fuente principal: RIGIDEZ_SIMULTANEA_DEL_LECTOR_RADIAL_20260926.md.
% Complementos demostrativos: COMPOSICION_METRICA_MEMORIA_Y_ACOPLAMIENTO.md
% y ACCION_ACOPLADA_FASE_Y_MEMORIA_DE_FRONTERA.md.
% La integración, la compilación y la revisión visual corresponden al editor.
\section{Simultaneous rigidity of the longitudinal reader and radial coupling}
\label{g26:sec:rigidez-radial-es}

The Triadic Modular Holography construction transmits to the radial realization an enriched state with its coordinates, orientation, carry, memory and incidences. The sequence
\[
 \begin{gathered}
 \mathrm{APP}\longrightarrow\mathrm{TRIT}\longrightarrow\mathrm{TPK}
 \longrightarrow\text{enriched state}\\
 \longrightarrow\text{joint discrete structure of the continuum}
 \end{gathered}
\]
preserves the Archimedean projection and incidence projection, together with the five inherent constructions of the continuum. The extension
\[
 w_6\longrightarrow w_{12}\longrightarrow w_{18}\longrightarrow
 w_{24}\longrightarrow w_{30}\longrightarrow R_{36}\longrightarrow G_9
\]
is articulated with the operator order \(U_t\longrightarrow\Gamma_9\longrightarrow K_{\rm ph}\): phase returns and the memory record increases. The regional coordinates, selected record \(K\), coordinate \(\alpha\) and return action are received with this common provenance and with their previously constructed readers.

The result of this section is a classification of realizations simultaneously preserving the generated data and the longitudinal, metric and temporal rules specified below. For each positive character the radial response is determined; the coefficient relating that response to mass is common to all admissible characters. The gravitational interpretation is formulated after both readings have been constructed on the same space.

\subsection{Incidence record and constitutive constraints}

Consider complex Hilbert spaces \(V_{\rm in}\) and \(V_{\rm out}\) of dimension \(N\) with fixed inner products. The marked cardinality specified in R1 is positive; in particular, both fibres are nonzero. Adjoints are taken with respect to these products. The symbol \(\mathsf U\) will denote normalized archive transport, and \(U_t\) will retain its meaning as evolution of the enriched state.

\paragraph{R1. Marked record.}
The incidence set is
\begin{equation}
 \Omega=\left\{(j,k,q,u,v):
 \begin{array}{l}
 j,k\in\mathbb Z/12\mathbb Z,\quad k-j\in\{0,3,6,9\},\\
 1\leq q\leq8,\quad u<v,\quad
 u,v\in\{1,2,4,5,7,8\}
 \end{array}\right\}.
 \label{g26:eq:omega-rigidez-es}
\end{equation}
The twelve phases, four relative positions, eight positions of index \(q\) and fifteen pairs in the last set give
\begin{equation}
 N=|\Omega|=12\cdot4\cdot8\binom62
   =12\cdot4\cdot120=5760.
 \label{g26:eq:cardinal-rigidez-es}
\end{equation}
Each incidence has its individual mark. In the corresponding basis \((e_i)_{i\in\Omega}\) of \(V_{\rm out}\), let \(P_i=e_ie_i^*\). The common mode \(\boldsymbol u\) is normalized by \(|u_i|^2=1/N\); write \(P=\boldsymbol u\boldsymbol u^*\). The marked algebra is \(\mathcal A_\Omega=\operatorname{span}\{P_i:i\in\Omega\}\).

\paragraph{R2. Normalization and complete archive.}
The generated inverter \(B:V_{\rm in}\to V_{\rm out}\) satisfies
\begin{equation}
 B^*B=\tfrac14 I_{\rm in},\qquad
 BB^*=\tfrac14 I_{\rm out}.
 \label{g26:eq:normalizacion-b-es}
\end{equation}
The two components are preserved
\begin{equation}
 C_0=(I-P)B,\qquad M_0=PB,\qquad
 C_0+M_0=B,\qquad \mathsf U=2B.
 \label{g26:eq:archivo-completo-es}
\end{equation}
The first records variation relative to the common mode; the second preserves that mode. Their images are orthogonal and their sum recovers the complete transport. By \eqref{g26:eq:normalizacion-b-es}, \(\mathsf U\) is unitary.

\paragraph{R3. Individual inscription.}
The map \(\mathcal E:\operatorname{End}(V_{\rm out})\to\mathcal A_\Omega\) is a linear conditional expectation fixing each \(P_i\) and satisfying
\begin{equation}
 \mathcal E(A_1TA_2)=A_1\mathcal E(T)A_2,
 \qquad A_1,A_2\in\mathcal A_\Omega.
 \label{g26:eq:bimodularidad-es}
\end{equation}
Inscription of the result preserves the complete construction data in a complementary record. The expectation determines the inscribed observable; archive preservation determines its subsequent reconstruction.

\paragraph{R4. Longitudinal composition.}
Fix the generated coordinate \(\alpha>0\) and the longitudinal reference section \(L_*>0\). Longitudinal transport is defined by
\begin{equation}
 D=L_*\alpha^{16}\mathcal E(C_0C_0^*)\mathsf U:
 V_{\rm in}\longrightarrow V_{\rm out}.
 \label{g26:eq:regla-longitudinal-es}
\end{equation}
This rule applies the inscribed coefficient linearly to the longitudinal section. Its type fixes the position of the power \(\alpha^{16}\) and return coefficient within the composition. Reading an intensity root is a different operation with its own transformation rule.

\paragraph{R5. Positive radial metric.}
Let \(X=X^*>0\) be the complete character on \(V_{\rm in}\). Define the positive radial reading \(R_X\) on \(V_{\rm out}\) by
\begin{equation}
 \|R_Xv\|^2=\|XD^*v\|^2
 \qquad(v\in V_{\rm out}).
 \label{g26:eq:metrica-radial-es}
\end{equation}
This equality fixes the metric on every vector of the fibre. In polar notation, \(R_X\) is the arrival modulus \(\lvert(DX)^*\rvert=\sqrt{DX^2D^*}\).

\paragraph{R6. Action and clock.}
The generated action section \(\hbar_{\rm ret}>0\) and the real phase lift satisfy
\begin{equation}
 S(\theta)=\hbar_{\rm ret}\theta.
 \label{g26:eq:accion-fase-es}
\end{equation}
The constitutive velocity \(c>0\) is preserved. Once the intensive length \(L\) of \eqref{g26:eq:regla-longitudinal-es} is obtained, the conjugate clock satisfies \(\omega L=c\). For \(\theta(t)=\omega t\), the action rate determines the energy scale
\begin{equation}
 E_{L}=\frac{dS}{dt}=\hbar_{\rm ret}\omega
   =\frac{\hbar_{\rm ret}c}{L}.
 \label{g26:eq:escala-energetica-es}
\end{equation}
The lift and its refinements remain associated with route memory.

\paragraph{R7. Readings of the same character.}
On \(Y=\mathsf U X\mathsf U^*>0\), define
\begin{equation}
 H_X=E_{L}Y,\qquad M_X=\frac{H_X}{c^2},\qquad
 \Lambda_X=LY^{-1}.
 \label{g26:eq:lecturas-conjugadas-es}
\end{equation}
Energy, mass and conjugate length preserve the same character and transport. The inverse acts in the positive sector; the null sector retains its own treatment.

\paragraph{R8. Reduced radius and dimensional sections.}
The physical identification of this realization is
\begin{equation}
 R_X=\frac{G}{c^2}M_X.
 \label{g26:eq:radio-reducido-es}
\end{equation}
Thus \(R_X\) represents the reduced gravitational radius. The corresponding Schwarzschild radius is \(2R_X\). \(L_*\), the action section and the constitutive velocity section are jointly maintained; a change of units transforms their coordinates according to their dimensions.

\paragraph{Provenance of the constraints.}
Hadamard transport, archive preservation, inscription, phase action and reciprocity have antecedents in the corpus. The expanded record \eqref{g26:eq:omega-rigidez-es} and compositions \eqref{g26:eq:regla-longitudinal-es}--\eqref{g26:eq:metrica-radial-es} are explicitly incorporated in the assembled construction. R4 and R5 are constitutive rules of this realization; the following proof establishes their joint consequence with the remaining antecedents. R8 specifies the physical identification fixing the coefficient's meaning.

\subsection{Simultaneous determination theorem}

\paragraph{Theorem.}
For the same generated data \(B\), \((P_i)\), \(P\), \(\alpha\), \(L_*\), \(\hbar_{\rm ret}\), \(c\) and the same character \(X\), constraints R1--R8 uniquely determine inscription, longitudinal transport and radial response. For all admissible positive characters, the coefficient \(G\) is common. One has
\begin{align}
 \mathcal E&=\Delta_\Omega,
 &\Delta_\Omega(T)&=\sum_{i\in\Omega}P_iTP_i,\label{g26:eq:delta-unica-es}\\
 r_N&=\frac{N-1}{4N}=\frac{5759}{23040},
 &L&=L_*\alpha^{16}r_N,\label{g26:eq:longitud-unica-es}\\
 D&=L\mathsf U,
 &R_X&=L\mathsf U X\mathsf U^*,\label{g26:eq:radial-unica-es}\\
 R_X\Lambda_X&=L^2I,
 &H_X\Lambda_X&=\hbar_{\rm ret}cI.\label{g26:eq:productos-radiales-es}
\end{align}
The coupling is determined by
\begin{equation}
 \boxed{G=\frac{c^3L^2}{\hbar_{\rm ret}}
 =\frac{c^3L_*^2}{\hbar_{\rm ret}}
   \left(\frac{5759}{23040}\right)^2\alpha^{32}.}
 \label{g26:eq:g-rigido-es}
\end{equation}

\paragraph{Proof: uniqueness of inscription.}
Let \(E_{ij}=e_ie_j^*\) be the matrix units. By bimodularity,
\[
 \mathcal E(E_{ij})=P_i\mathcal E(E_{ij})P_j.
\]
For \(i\ne j\), membership of \(\mathcal E(E_{ij})\) in the diagonal algebra implies \(P_i\mathcal E(E_{ij})P_j=0\). For \(i=j\), fixing \(P_i\) gives \(\mathcal E(E_{ii})=P_i\). Linearity then determines \eqref{g26:eq:delta-unica-es} on all matrices. Marked inscription is thus fixed by its algebraic properties.

\paragraph{Proof: return coefficient and archive.}
Normalization of \(B\) and \(P^2=P=P^*\) give
\[
 C_0C_0^*=\tfrac14(I-P),\qquad M_0M_0^*=\tfrac14P.
\]
Since \(P_iPP_i=|u_i|^2P_i=P_i/N\), one obtains
\begin{equation}
 \Delta_\Omega(C_0C_0^*)=\frac{N-1}{4N}I,
 \qquad
 \Delta_\Omega(M_0M_0^*)=\frac{1}{4N}I.
 \label{g26:eq:retorno-y-memoria-es}
\end{equation}
Both readings sum to \(I/4\). The first is scalar, so any normalized distribution \((p_i)\) on incidences, with \(p_i\geq0\) and \(\sum_i p_i=1\), produces \(\sum_i p_i r_N=r_N\). The same coefficient is therefore preserved under arbitrary diagonal weighting.

\paragraph{Proof: length and metric.}
Substituting \eqref{g26:eq:retorno-y-memoria-es} into R4 gives \(D=L\mathsf U\), with \(L\) as in \eqref{g26:eq:longitud-unica-es}. Unitarity of \(\mathsf U\) implies \(DD^*=L^2I\). By R5, the quadratic forms of \(R_X^2\) and \(DX^2D^*\) coincide. The polarization identity determines
\[
 R_X^2=DX^2D^*=L^2\mathsf U X^2\mathsf U^*.
\]
The operator \(L\mathsf U X\mathsf U^*\) is positive and its square is the operator on the right-hand side. Uniqueness of the positive square root proves \eqref{g26:eq:radial-unica-es}. In particular, \(R_I=LI\).

\paragraph{Proof: action, reciprocity and coupling.}
R6 and R7 give
\[
 H_X=\frac{\hbar_{\rm ret}c}{L}Y,
 \qquad M_X=\frac{\hbar_{\rm ret}}{cL}Y,
 \qquad \Lambda_X=LY^{-1}.
\]
Multiplication produces both products of \eqref{g26:eq:productos-radiales-es}. Next, R8 is equivalent to
\[
 LY=\frac{G\hbar_{\rm ret}}{c^3L}Y.
\]
Invertibility of \(Y\) determines \eqref{g26:eq:g-rigido-es}. The same substitution verifies the equality for each positive character. If \(G_1\) and \(G_2\) satisfy the identity with the same data, \((G_1-G_2)Y=0\), and therefore \(G_1=G_2\). This conclusion distinguishes the state-dependent response \(R_X\) from the common coefficient \(G\). \hfill\(\square\)

\subsection{Reciprocity and phase rigidity}

The power--logarithmic antecedent uses, for \(x>0\),
\[
 \ell_p(x)=\frac{x^p-1}{p}\quad(p\ne0),
 \qquad \ell_0(x)=\log x.
\]
Direct substitution proves \(\ell_{-p}(x^{-1})=-\ell_p(x)\). On the positive domain of their inverses, this identity is equivalent to
\[
 \exp_p(s)\exp_{-p}(-s)=1.
\]
Spectral composition on a positive character transmits this product to the pair \(LY\), \(LY^{-1}\). In the present realization, R5 can equivalently be expressed, once R1--R4 and R7 are fixed, by \(R_X\Lambda_X=L^2I\): multiplication by \(\Lambda_X^{-1}\) determines \(R_X=L\mathsf U X\mathsf U^*\), and this expression satisfies R5.

The two products of \eqref{g26:eq:productos-radiales-es} also give a direct form of rigidity:
\begin{equation}
 R_X=\frac{L^2}{\hbar_{\rm ret}c}H_X.
 \label{g26:eq:proporcion-operadores-es}
\end{equation}
For a positive rescaling \(s\), the pair \((sR_X,\Lambda_X/s)\) preserves the radial product; the product \(H_X(\Lambda_X/s)=\hbar_{\rm ret}cI/s\), with energy and action fixed, requires \(s=1\). Joint preservation determines the scale. Likewise, any positive radial operator with the same metric R5 equals \(R_X\) by uniqueness of the positive root.

The lifted action is also exactly determined. For \(\theta\ne0\), the equality \(S/\hbar_{\rm ret}=\theta\) fixes \(S\). Even when rotations of all subdivisions \(\theta/9^n\) are preserved, a fixed real rescaling \(a\) satisfying
\[
 \exp(-ia\theta/9^n)=\exp(-i\theta/9^n)
 \qquad(n\geq0)
\]
satisfies \((a-1)\theta/9^n\in2\pi\mathbb Z\). This sequence tends to zero; for sufficiently large \(n\) its only possible value is zero. From \(\theta\ne0\) one deduces \(a=1\).

\subsection{Planck scales and specialization of the unit character}

Write \(L_\alpha=L\); on the return section, \(E_{P,\rm ret}=E_L\).

The constructed length $L_\alpha>0$, action section $\hbar_a>0$ and velocity $c>0$ determine the scales

\[
t_P=\frac{L_\alpha}{c},\qquad
m_{P,a}=\frac{\hbar_a}{cL_\alpha},\qquad
E_{P,a}=\frac{\hbar_a c}{L_\alpha}.
\]

Indeed, substituting $G_a=c^3L_\alpha^2/\hbar_a$ into the expressions $\sqrt{\hbar_aG_a/c^5}$, $\sqrt{\hbar_ac/G_a}$ and $\sqrt{\hbar_ac^5/G_a}$ gives those three identities respectively. Positivity fixes the root in each case. Length precedes recognition of these scales; their conventional expressions are identities of the already constructed realization.

For a mode of character $y>0$, the mass and longitudinal readings are

\[
m(y)=m_{P,a}y,\qquad r_g(y)=L_\alpha y,\qquad
\bar\lambda(y)=\frac{L_\alpha}{y}.
\]

Thus $r_g(y)\bar\lambda(y)=L_\alpha^2$ for all positive modes. Equality of the two lengths characterizes the mode $y=1$, since $y=y^{-1}$ and $y>0$ imply $y=1$. In this mode, mass is $m_{P,a}$ and both lengths equal $L_\alpha$. General reciprocity is thus distinguished from its specialization to the unit character.

The chronological step and Planck time satisfy $t_0=(\pi_{\rm HMT}/54)t_P$. The complete 108-step phase period is $108t_0=2\pi_{\rm HMT}t_P$ and its angular frequency is $\omega_{108}=1/t_P=c/L_\alpha$. A horizon of radius $R_h=\zeta_hL_\alpha$, with $\zeta_h>0$, preserves its own geometric variable. The Schwarzschild radius associated with a mass is $2r_g$.

If the two action sections are compared with $L_\alpha$ and $c$ fixed, the ratio $\rho=\hbar_b/\hbar_a$ gives $G_b=G_a/\rho$, $m_{P,b}=\rho m_{P,a}$ and $E_{P,b}=\rho E_{P,a}$; $L_\alpha$ and $t_P$ remain equal. This comparison between sections remains distinct from temporal evolution of the constants.

\subsection{Archive covariance and transformation of units}

Let \(T:V_{\rm in}\to V'_{\rm in}\) and \(S:V_{\rm out}\to V'_{\rm out}\) be unitary. Joint transport of the data is
\[
 B'=SBT^*,\quad P'=SPS^*,\quad P_i'=SP_iS^*,\quad X'=TXT^*.
\]
The transported expectation \(\mathcal E'(A')=S\mathcal E(S^*A'S)S^*\) fixes \(P_i'\) and is bimodular in the transformed marked algebra. By substitution,
\begin{align*}
 \mathsf U'&=S\mathsf UT^*,&D'&=SDT^*,\\
 R'_{X'}&=SR_XS^*,&H'_{X'}&=SH_XS^*,&
 \Lambda'_{X'}&=S\Lambda_XS^*.
\end{align*}
Conjugation preserves the R1--R8 equalities and the same coefficient \(G\). Marks, orientations and records are transported with the bases. Incidence permutations are a particular case.

Incorporating an auxiliary memory space \(\mathcal F\) transports \(P\) to \(P\otimes I_{\mathcal F}\), \(\mathsf U\) to \(\mathsf U\otimes I_{\mathcal F}\) and \(\Delta_\Omega\) to \(\Delta_\Omega\otimes\operatorname{id}\). The inscription of \(C_0C_0^*\otimes I_{\mathcal F}\) is \(r_N I\otimes I_{\mathcal F}\); the coefficient preserves the original marked cardinality.

To specify dimensional covariance, let \(\ell\), \(\tau\) and \(m\) be the positive factors of the new length, time and mass units relative to the previous units. Numerical coordinates satisfy
\[
 L'=L/\ell,\qquad c'=c\tau/\ell,\qquad
 \hbar'_{\rm ret}=\hbar_{\rm ret}\tau/(m\ell^2).
\]
Thus,
\[
 \frac{(c')^3(L')^2}{\hbar'_{\rm ret}}
 =G\frac{m\tau^2}{\ell^3},
\]
which is the transformation of a quantity of dimension \(\mathrm{longitud}^3\mathrm{masa}^{-1}\mathrm{tiempo}^{-2}\). The section \(L_*\) transforms as \(L\), while \(\alpha\) and \(r_N\) remain dimensionless.

\subsection{Exact elimination of coupled memory}

Let \(\mathcal K\) and \(\mathcal H\) be finite Hilbert spaces, \(W=W^*>0\) a global weight, and \(\mathcal C:\mathcal K\to\mathcal H\) surjective. The action of residuals \(r\), with boundary defect \(z\), is
\[
 \mathcal S(r)=\mathfrak a\,r^*Wr,
 \qquad \mathcal Cr=z,\qquad \mathfrak a>0.
\]
The global weight preserves its cross blocks between incidences. Define \(\mathcal R=\mathcal CW^{-1}\mathcal C^*\). For \(v\ne0\), injectivity of \(\mathcal C^*\) gives
\[
 v^*\mathcal Rv=\|W^{-1/2}\mathcal C^*v\|^2>0.
\]
Thus \(\mathcal R\) is invertible. The least-action residual is
\[
 r_*(z)=W^{-1}\mathcal C^*\mathcal R^{-1}z.
\]
One verifies \(\mathcal Cr_*=z\). Each admissible residual is written \(r=r_*+k\), with \(\mathcal Ck=0\) and \(k^*Wr_*=k^*\mathcal C^*\mathcal R^{-1}z=0\). Consequently,
\begin{equation}
 \mathcal S(r)=\mathfrak a\,z^*\mathcal R^{-1}z
             +\mathfrak a\,k^*Wk.
 \label{g26:eq:accion-residuo-es}
\end{equation}
Positivity of \(W\) proves uniqueness of the minimum. The pair \((z,k)\) preserves and reconstructs the complete residual.

For a route with transports \(U_j\), the residuals \(r_j=f_j-U_jf_{j-1}\) are reconstructed through \(f_j=U_jf_{j-1}+r_j\). The map \(\mathcal C\) collects the subsequent transports of each residual towards the endpoint. If its blocks are \(\mathcal C_j\), then
\[
 \mathcal R=\sum_{j,k}\mathcal C_j(W^{-1})_{jk}\mathcal C_k^*.
\]
This expression preserves correlations between route segments. For an invertible longitudinal transport \(D\), the boundary \(y=Dz\) has response
\[
 \mathcal R_L=D\mathcal RD^*,\qquad
 \mathcal S_{{\rm ef},L}(y)
 =\mathfrak a\,y^*(D\mathcal RD^*)^{-1}y
 =\mathfrak a\,z^*\mathcal R^{-1}z.
\]
The same action is obtained whether memory is eliminated first or the boundary is transported first. The prefactor preserves the normalization of the source action; when that action carries \(\mathfrak a=\hbar_{\rm ret}\), reduction transmits the same section.

\subsection{Joint reduction of radius, energy and conjugate length}

Write the same positive character in a boundary and memory decomposition:
\[
 Y=\begin{pmatrix}A&F\\F^*&C\end{pmatrix},\qquad C>0,
 \qquad Y_{\rm ef}=A-FC^{-1}F^*.
\]
The identity
\[
 \begin{pmatrix}b\\y\end{pmatrix}^{\!*}Y
 \begin{pmatrix}b\\y\end{pmatrix}
 =b^*Y_{\rm ef}b+\eta^*C\eta,
 \qquad \eta=y+C^{-1}F^*b,
\]
proves \(Y_{\rm ef}>0\), characterizes the least-energy extension and preserves the reconstruction \(y=\eta-C^{-1}F^*b\). For \(s>0\), substitution into the Schur complement gives \(\operatorname{Schur}(sY)=sY_{\rm ef}\). Consequently,
\[
 R_{\rm ef}=LY_{\rm ef},\qquad
 H_{\rm ef}=E_{L}Y_{\rm ef}.
\]
To calculate the compressed conjugate length, solve \(Y(x,y)=(f,0)\). The second equation gives \(y=-C^{-1}F^*x\), and the first gives \(x=Y_{\rm ef}^{-1}f\). The boundary block of \(Y^{-1}\) is therefore \(Y_{\rm ef}^{-1}\). Thus,
\begin{equation}
 \Lambda_{\rm b}=LY_{\rm ef}^{-1},\qquad
 R_{\rm ef}\Lambda_{\rm b}=L^2I,\qquad
 R_{\rm ef}=\frac{L}{E_{L}}H_{\rm ef}.
 \label{g26:eq:schur-reciproco-es}
\end{equation}
Reduction acts on the same blocks in both readings and preserves the residual \(\eta\). The effective character may depend on memory in matrix form, while the coefficient \(L/E_{L}\) remains fixed.

\subsection{Dynamic memory and induced temporal norm}

Dynamic elimination preserves an additional spectral dependence. In a stationary representation, let
\[
 H=\begin{pmatrix}H_{\rm bb}&V\\V^*&H_{\rm ii}\end{pmatrix},
 \qquad H_{\rm ii}>0.
\]
For \(z\) in the relevant resolvent domain, elimination of the interior component in \((H-zI)(b,y)=(f,0)\) gives
\begin{equation}
 \mathcal K_H(z)=H_{\rm bb}-zI
       -V(H_{\rm ii}-zI)^{-1}V^*.
 \label{g26:eq:resolvente-memoria-es}
\end{equation}
When both inverses exist, the boundary block of \((H-zI)^{-1}\) is \(\mathcal K_H(z)^{-1}\). For \(|z|<\lambda_{\min}(H_{\rm ii})\), the convergent resolvent series gives
\begin{align}
 \mathcal K_H(z)&=H_{\rm est}-zZ-z^2VH_{\rm ii}^{-3}V^*-\cdots,
 \label{g26:eq:expansion-resolvente-es}\\
 H_{\rm est}&=H_{\rm bb}-VH_{\rm ii}^{-1}V^*,\qquad
 Z=I+VH_{\rm ii}^{-2}V^*\geq I.
 \label{g26:eq:norma-z-es}
\end{align}
Positivity follows from \(Z-I=(H_{\rm ii}^{-1}V^*)^*(H_{\rm ii}^{-1}V^*)\). The static extension \((b,-H_{\rm ii}^{-1}V^*b)\) has squared norm \(b^*Zb\); energy and temporal norm arise from the same extension. The physical parameter is \(z=\hbar_{\rm ret}\omega\).

To first temporal order, \(T=Z^{1/2}\) and \(\psi=Tb\) transform the pair \((H_{\rm est},Z)\) into \((T^{-*}H_{\rm est}T^{-1},I)\). The section \(\hbar_{\rm ret}\) remains fixed. If \(T\) depends on time, the transformation also preserves the term \(\dot T b\) of \(\dot\psi=\dot T b+T\dot b\).

Radial proportionality is preserved in the complete resolvent. Writing \(\chi_{\rm rad}=L/E_{L}\) and \(R=\chi_{\rm rad} H\), substitution of each block proves
\begin{equation}
 \mathcal K_R(\chi_{\rm rad} z)=\chi_{\rm rad}\mathcal K_H(z).
 \label{g26:eq:resolvente-proporcional-es}
\end{equation}
This identity transports the operator and spectral parameter simultaneously and preserves every order of dynamic memory. For the effective forms of \eqref{g26:eq:schur-reciproco-es}, canonical normalization requires dual transport of the inverse length:
\begin{align*}
 H_{\rm can}&=T^{-*}H_{\rm ef}T^{-1},&
 R_{\rm can}&=T^{-*}R_{\rm ef}T^{-1},&
 \Lambda_{\rm can}&=T\Lambda_{\rm b}T^*.
\end{align*}
Multiplying these expressions gives
\[
 R_{\rm can}=\frac{L}{E_{L}}H_{\rm can},\qquad
 R_{\rm can}\Lambda_{\rm can}=L^2I.
\]
The proof allows noncommuting \(T\), \(R_{\rm ef}\) and \(\Lambda_{\rm b}\). The induced temporal norm and dual length jointly preserve the same \(G\).

\subsection{Refinement, spectral domains and scope}

Let \(J\) and \(K\) be isometric inscriptions between two levels satisfying \(X'J=JX\) and \(\mathsf U'J=K\mathsf U\), with the same intensive section \(L\). Then
\[
 R'K=L\mathsf U'X'\mathsf U'^*K
     =KL\mathsf U X\mathsf U^*=KR,
\]
and the same calculation proves \(H'K=KH\). Intertwining of the inverses is formulated on the corresponding domains. When the compatible system has a positive injective self-adjoint limit \(X\), spectral calculus defines \(Y=\mathsf U X\mathsf U^*\), \(R=LY\), \(H=E_{L}Y\) and \(\Lambda=LY^{-1}\). Proportionality \(R=(L/E_{L})H\) holds on \(\operatorname{Dom}(Y)\), whereas \(R\Lambda=L^2I\) and \(H\Lambda=\hbar_{\rm ret}cI\) hold on \(\operatorname{Dom}(Y^{-1})\), with bounded extensions of their products. Spectral cuts \(\mathbf1_{[\varepsilon,M]}(Y)\) verify the identities before passage to the limit on these domains. A refined edge length \(L/9\) incorporates its own scale factor; the intensive section is preserved according to the declared intertwining.

The classification distinguishes transformations preserving the realization from those modifying a constitutive rule. Diagonal weightings, unitary equivalences, the auxiliary archive and joint reduction preserve \(G\). Changing the marked resolution modifies R1; replacing \(r_N\) with \(\sqrt{r_N}\), \(\alpha^{16}\) with its square, or the inscribed component with the total norm modifies R4. Isolated radius rescaling changes R5, and isolated action rescaling changes R6. R8 maintains the distinction between reduced radius, Schwarzschild radius and change of units.

One thus obtains mathematical uniqueness of the R1--R8 realization: the generated data and character fix the response, and the common composition fixes \(G\). The longitudinal and metric rules appear explicitly in this conclusion as incorporated compositions; gravitational identification retains its physical status. The proof covers the declared class with its constitutive definitions and compatible transformations. The rational checks of the development verify finite cases of the stated identities; the general proof resides in the preceding equations and arguments.
% END CR28_RADIAL_FULL

